\documentclass[12pt]{article}
%^ Start of document
\usepackage{times}  %Makes text TimesNewRoman
\usepackage{setspace} %Allows you to set single or double space 
\usepackage{biblatex} %Imports biblatex package
\usepackage{graphicx} % Required for inserting images
\usepackage{authblk}
\usepackage[colorlinks=true, urlcolor=blue, linkcolor=blue]{hyperref}
\usepackage{subfig}
\usepackage{float}
\usepackage{tabularx}
\usepackage{multirow}
\usepackage[paperheight=11in,
   paperwidth=8.5in,
   top=20mm,
   bottom=20mm,
   left=40mm,
   right=40mm]{geometry}
\usepackage{moresize}
\usepackage{tikz}
\usepackage{xcolor}
\usepackage{physics}
\usepackage{amssymb}
\usepackage{mathrsfs}
\usepackage{amsmath}
\usepackage[document]{ragged2e}
\usepackage{comment}
\usepackage{xcolor}
\addbibresource{Qhw.bib}
\begin{document}
\singlespacing  %Sets single spacing for whole document
\title{Quantum CH 11 HW}

\author[1]{Zachary Tullier}
\affil[1]{Student\\Physics Department\\ University of Houston - Clear Lake \\Houston, Texas\\U.S}
\date{}
\maketitle

\section{Textbook Question 11.3}
    Consider the scattering from the potential \( V(r) = V_0 e^{-r^2/a^2} \). Find  
    \begin{itemize}
        \item [(a)] the differential cross section in the first Born approximation
        \item [(b)] the total cross section
    \end{itemize}

    \subsection{Solution, part (a):}
        In the first Born approximation, the scattering amplitude for a spherically symmetric potential is:
        \[
            f(\theta) = -\frac{2m}{\hbar^2} \int_0^\infty V(r) \frac{\sin(qr)}{q} r \, dr
        \]
        
        where 
        \[ 
            q = 2k \sin\left( \frac{\theta}{2} \right) 
        \]
        
        Substitute the given potential:
        \[
            f(\theta) = -\frac{2m V_0}{\hbar^2 q} \int_0^\infty r e^{-r^2/a^2} \sin(qr) \, dr
        \]
        
        This is a known integral:
        \[
            \int_0^\infty r e^{-r^2/a^2} \sin(qr) \, dr = \frac{q a^2}{2} e^{-q^2 a^2 / 4}
        \]
        
        Thus:
        \[
            f(\theta) = -\frac{2m V_0}{\hbar^2 q} \cdot \frac{q a^2}{2} e^{-q^2 a^2 / 4} = -\frac{m V_0 a^2}{\hbar^2} e^{-q^2 a^2 / 4}
        \]
        
        The differential cross section is:
        \[
            \boxed{
            \frac{d\sigma}{d\Omega} = \left| f(\theta) \right|^2 = \left( \frac{m V_0 a^2}{\hbar^2} \right)^2 \exp\left[ -\frac{q^2 a^2}{2} \right]}
        \]
        
        with 
        \[
            q = 2k \sin\left( \frac{\theta}{2} \right) 
        \]

    \subsection{Solution, part (b):}
        The total cross section is obtained by integrating over all solid angles:
        \[
            \sigma = \int \frac{d\sigma}{d\Omega} \, d\Omega = 2\pi \int_0^\pi \left| f(\theta) \right|^2 \sin\theta \, d\theta
        \]
        
        Substitute the expression for \( f(\theta) \):
        \[
            \sigma = 2\pi \left( \frac{m V_0 a^2}{\hbar^2} \right)^2 \int_0^\pi \exp\left[ -4k^2 a^2 \sin^2\left( \frac{\theta}{2} \right) \right] \sin\theta \, d\theta
        \]
        
        Use u-substitution \( \theta' = \theta/2 \), then \( \sin\theta = 2\sin\theta' \cos\theta' \), \( d\theta = 2 d\theta' \)
        \[
            \sigma = 8\pi \left( \frac{m V_0 a^2}{\hbar^2} \right)^2 \int_0^{\pi/2} \sin\theta' \cos\theta' \, e^{-4k^2 a^2 \sin^2\theta'} \, d\theta'
        \]
        
        Now substitute \( u = \sin\theta' \Rightarrow du = \cos\theta' d\theta' \), limits: \( u = 0 \to 1 \)
        \[
            \sigma = 8\pi \left( \frac{m V_0 a^2}{\hbar^2} \right)^2 \int_0^1 u e^{-4k^2 a^2 u^2} \, du
        \]
        
        This integral evaluates to:
        \[
            \int_0^1 u e^{-bu^2} \, du = \frac{1 - e^{-b}}{2b}, \quad \text{with } b = 4k^2 a^2
        \]
        
        So:
        \[
            \sigma = 8\pi \left( \frac{m V_0 a^2}{\hbar^2} \right)^2 \cdot \frac{1 - e^{-4k^2 a^2}}{8k^2 a^2}
            = \pi \left( \frac{m V_0 a}{\hbar^2 k} \right)^2 \left( 1 - e^{-4k^2 a^2} \right)
        \]
        
        \[
            \boxed{\sigma = \pi \left( \frac{m V_0 a}{\hbar^2 k} \right)^2 \left( 1 - e^{-4k^2 a^2} \right)}
        \]


\section{Textbook Question 11.5}
    Consider the elastic scattering from the delta potential \( V(r) = V_0 \delta(r - a) \).
    
    \begin{enumerate}
        \item[(a)] Calculate the differential cross section in the first Born approximation.
        \item[(b)] Find an expression between \( V_0, a, \mu, \) and \( k \) so the Born approximation is valid.
    \end{enumerate}

    \subsection{Solution, part (a):}
        The first Born approximation for a spherically symmetric potential gives the scattering amplitude:
        \[
            f(\theta) = -\frac{2\mu}{\hbar^2 q} \int_0^\infty V(r)\, r\, \sin(qr) \, dr
        \]
        
        where \( q = 2k \sin\left( \frac{\theta}{2} \right) \), \( \mu \) is the reduced mass, and \( k = \sqrt{2\mu E} / \hbar \).
    
        Substituting the delta-shell potential \( V(r) = V_0 \delta(r - a) \):
        \[
            f(\theta) = -\frac{2\mu V_0}{\hbar^2 q} \cdot a \sin(qa)
        \]
    
        Therefore, the differential cross section is:
        \[
            \frac{d\sigma}{d\Omega} = |f(\theta)|^2 = \left( \frac{2\mu a V_0}{\hbar^2 q} \right)^2 \sin^2(qa)
        \]
    
        Substitute \( q = 2k \sin\left( \frac{\theta}{2} \right) \) and \( qa = 2ka \sin\left( \frac{\theta}{2} \right) \):
        \[
            \boxed{
            \frac{d\sigma}{d\Omega} = \left( \frac{\mu a V_0}{\hbar^2 k \sin(\theta/2)} \right)^2 \sin^2\left( 2ka \sin\left( \frac{\theta}{2} \right) \right)}
        \]

    \subsection{Solution, part (b):}
        The Born approximation is valid when the scattering amplitude is small:
        \[
            |f(\theta)| \ll 1
        \]
    
        From part (a), the maximum of \( \sin(qa) \leq 1 \), so:
        \[
            |f(\theta)|_{\max} = \frac{2\mu a V_0}{\hbar^2 q}
        \]
    
        To ensure validity for all \( \theta \), we impose:
        \[
            \boxed{
            \frac{2\mu a V_0}{\hbar^2} \ll 1            }
        \]
    
        Or more conservatively, accounting for typical \( q \sim k \):
        \[
            \boxed{
            \frac{2\mu a V_0}{\hbar^2 k} \ll 1
        }
        \]


\section{Textbook Question 11.7}
    Find the differential cross section in the first Born approximation for the elastic scattering of a particle of mass \( m \), which is initially traveling along the \( z \)-axis, from a nonspherical, double-delta potential:
    \[
        V(\vec{r}) = V_0 \delta(\vec{r} - a \vec{k}) - V_0 \delta(\vec{r} + a \vec{k}),
    \]
    
    where \( \vec{k} \) is the unit vector along the \( z \)-axis.

    \subsection{Solution:}
        We are given a non-spherically symmetric double-delta potential:
        \[
            V(\vec{r}) = V_0 \delta(\vec{r} - a \hat{k}) - V_0 \delta(\vec{r} + a \hat{k}),
        \]
        where \( \hat{k} \) is the unit vector along the \( z \)-axis.
        
        The scattering amplitude in the first Born approximation is:
        \[
            f(\vec{q}) = -\frac{2\mu}{\hbar^2} \frac{1}{4\pi} \int e^{i \vec{q} \cdot \vec{r}} V(\vec{r}) \, d^3r
        \]
        
        Substituting the potential:
        \[
            f(\vec{q}) = -\frac{2\mu V_0}{\hbar^2} \frac{1}{4\pi} \left[ e^{i q_z a} - e^{-i q_z a} \right]
            = -i \frac{\mu V_0}{\pi \hbar^2} \sin(q_z a)
        \]
        
        Differential Cross Section
        \[
            \frac{d\sigma}{d\Omega} = |f(\vec{q})|^2 = \left( \frac{\mu V_0}{\pi \hbar^2} \right)^2 \sin^2(q_z a)
        \]
        
        Since the incident direction is along the \( z \)-axis and \( \vec{q} = \vec{k} - \vec{k}' \), we get:
        \[
            q_z = k \cos\theta - k = -2k \sin^2\left( \frac{\theta}{2} \right)
        \]
        
        Substitute into the result:
        \[
            \boxed{
            \frac{d\sigma}{d\Omega} = \left( \frac{\mu V_0}{\pi \hbar^2} \right)^2 \sin^2\left( 2ka \sin^2\left( \frac{\theta}{2} \right) \right)        }
        \]

\section{Textbook Question 11.9}
    Consider the elastic scattering of a particle of mass \( m \) and initial momentum \( \hbar k \) off a delta potential 
    \[
        V(\vec{r}) = V_0 \delta(x) \delta(y) \delta(z - a),
    \]
    where \( V_0 \) is a constant.
    
    \begin{enumerate}
        \item[(a)] What is the physical dimension of the constant \( V_0 \)?
        \item[(b)] Calculate the differential cross section in the first Born approximation.
        \item[(c)] Repeat (b) for the case where the potential is now given by
        \[
        V(\vec{r}) = V_0 \delta(x) \left[ \delta(y - b)\delta(z) + \delta(y)\delta(z - a) \right].
        \]
    \end{enumerate}

    \subsection{Solution, part (a):}
    Since the delta functions have dimensions of inverse length:
    \[
        [V(\vec{r})] = \text{energy}, \quad [\delta(x)\delta(y)\delta(z)] = L^{-3} \Rightarrow [V_0] = \text{energy} \cdot \text{volume}
    \]
    \[
        \boxed{[V_0] = \text{Joule} \cdot \text{m}^3}
    \]

    \subsection{Solution, part (b):}
    Compute differential cross section (First Born Approximation). The scattering amplitude is:
    \[
        f(\vec{q}) = -\frac{2\mu}{\hbar^2} \frac{1}{4\pi} \int e^{i \vec{q} \cdot \vec{r}} V(\vec{r}) \, d^3r
        = -\frac{2\mu V_0}{\hbar^2} \frac{1}{4\pi} e^{i q_z a}
    \]

    The differential cross section is:
    \[
        \boxed{
        \frac{d\sigma}{d\Omega} = \left( \frac{\mu V_0}{2\pi \hbar^2} \right)^2    }
    \]

    \subsection{Solution, part (c):}
        Find the new potential:
        \[
            V(\vec{r}) = V_0 \delta(x)\left[ \delta(y - b)\delta(z) + \delta(y)\delta(z - a) \right]
        \]
    
        Then:
        \[
            f(\vec{q}) = -\frac{2\mu V_0}{\hbar^2} \frac{1}{4\pi} \left( e^{i q_y b} + e^{i q_z a} \right)
        \]
    
        So:
        \[
            \frac{d\sigma}{d\Omega} = \left( \frac{\mu V_0}{2\pi \hbar^2} \right)^2 \left| e^{i q_y b} + e^{i q_z a} \right|^2
            = 2 \left( \frac{\mu V_0}{2\pi \hbar^2} \right)^2 \left[ 1 + \cos(q_y b - q_z a) \right]
        \]
    
        Use \( q_y = k \sin\theta \sin\phi \), \( q_z = k(1 - \cos\theta) \):
    
        \[
            \boxed{
            \frac{d\sigma}{d\Omega} = 2 \left( \frac{\mu V_0}{2\pi \hbar^2} \right)^2 \left[ 1 + \cos\left( k b \sin\theta \sin\phi - k a (1 - \cos\theta) \right) \right]        }
        \]


\section{Textbook Question 11.12}
    Consider the \( S \)-partial wave scattering (\( l = 0 \)) between two identical spin \( 1/2 \) particles where the interaction potential is given approximately by
    \[
        \hat{V}(r) = V_0 \vec{S}_1 \cdot \vec{S}_2 \delta(r - a),
    \]
    where \( \vec{S}_1 \) and \( \vec{S}_2 \) are the spin vector operators of the two particles, and \( V_0 > 0 \). Assuming that the incident and target particles are unpolarized, find the differential and total cross sections.

    \subsection{Solution:}
        For two spin-1/2 particles, the possible total spin states are:
        \begin{itemize}
            \item Singlet: \( S = 0 \), one state
            \item Triplet: \( S = 1 \), three states
        \end{itemize}
    
        The operator \( \vec{S}_1 \cdot \vec{S}_2 \) acts as:
        \[
            \vec{S}_1 \cdot \vec{S}_2 = \frac{1}{2} \left[ S(S+1) - \frac{3}{2} \right] \hbar^2
        \]
    
        So:
        \[
            \vec{S}_1 \cdot \vec{S}_2 =
            \begin{cases}
                -\frac{3}{4} \hbar^2 & \text{for } S = 0 \ (\text{singlet}) \\
                \ \frac{1}{4} \hbar^2 & \text{for } S = 1 \ (\text{triplet})
            \end{cases}
        \]
    
        For each spin channel, the S-wave cross section is:
        \[
            \sigma = \frac{4\pi}{k^2} \left( \frac{\mu \lambda}{\hbar^2 k} \cdot \frac{\sin(ka)}{ka} \right)^2
        \]
    
        where \( \lambda = V_0 \cdot (\vec{S}_1 \cdot \vec{S}_2) \). So:
    
        \[
            \sigma_s = \frac{4\pi}{k^2} \left( \frac{\mu V_0}{k} \cdot \frac{-3}{4} \cdot \frac{\sin(ka)}{ka} \right)^2,
            \quad
            \sigma_t = \frac{4\pi}{k^2} \left( \frac{\mu V_0}{k} \cdot \frac{1}{4} \cdot \frac{\sin(ka)}{ka} \right)^2
        \]
    
        Unpolarized Averaging:
        \[
            \sigma_{\text{total}} = \frac{1}{4} \sigma_s + \frac{3}{4} \sigma_t
        \]
    
        Compute:
    
        \[
            \sigma_{\text{total}} = \frac{4\pi}{k^2} \left( \frac{\mu V_0 \sin(ka)}{k a} \right)^2 \cdot \left( \frac{9}{64} + \frac{3}{64} \right)
            = \frac{3\pi}{4k^2} \left( \frac{\mu V_0 \sin(ka)}{k a} \right)^2
        \]
    
        \[
            \boxed{            \sigma_{\text{total}} = \frac{3\pi}{4k^2} \left( \frac{\mu V_0 \sin(ka)}{k a} \right)^2        }
        \]
    
        Differential Cross Section (isotropic):
        \[
            \boxed{        \frac{d\sigma}{d\Omega} = \frac{\sigma_{\text{total}}}{4\pi} = \frac{3}{16k^2} \left( \frac{\mu V_0 \sin(ka)}{k a} \right)^2        }
        \]
    
\end{document}
